\documentstyle[%


righttag%


,11pt%


,amssymb%


,russian%               remove in Eng.


%,izvru%                 Set izven% in Eng.


,izvru-ks%             for brief communications


]{amsart}





%


\newcommand{\im}{\operatorname{Im}}


\newcommand{\Ker}{\operatorname{Ker}}


\newcommand{\tts}[1]{}





%\hoffset=-15mm%                  For Eng.


\setcounter{page}{77}


\god{2000}


\nomer{3~(454)} %                        Remove in Eng.


%\nomer{Vol.\,44, No.\,3}%               Set in Eng.


%\str{\pageref{begin}--\pageref{end}}%  Set in Eng.


\udk{517.988}





\author{\it Ф.К.\,Ахмадишина}


% NB:   ^^ remove in Eng.


\title{Об асимптотической оптимальности\\


проекционных методов}





\date{}


%\pagestyle{myheadings}%         Set in Eng.


\pagestyle{plain} %              Remove in Eng.


\begin{document}


%\thispagestyle{plain}%          Set in Eng.


\maketitle


\label{begin}








{\bf 1.} В связи с имеющимися многочисленными результатами по приближенному


решению различных классов операторных уравнений в последнее время


возникла и развивается теория оптимизации вычислительных методов,


т.\,е. теория построения и исследования наиболее точных


методов решения задач. В данной статье придерживаемся подхода,


разработанного Б.Г.\,Габдулхаевым [1]. Этот подход был использован


многими авторами при обосновании оптимальности (в основном, по порядку)


прямых методов решения операторных уравнений.





Рассмотрим уравнения


\begin{gather}


Kx=y \quad (x\in X,\ \; y\in Y),\\


P_nKx_n=P_ny \quad (x_n\in X_n).


\end{gather}


Здесь $X$, $Y$ --- нормированные пространства над $\Bbb C$ или $\Bbb R$,


$K:X \rightarrow Y$ --- линейная биекция,


$n$ --- фиксированное натуральное число, $X_n$ и $Y_n$ --- $n$-мерные


подпространства в $X$ и $Y$ соответственно, $P_n$ --- линейная сюръекция


из $Y$ на $Y_n$. Уравнение (2) называется {\it проекционным методом}


решения уравнения (1), заданным оператором $P_n$ и подпространством


приближенных решений $X_n$. Решение уравнения (1) обозначим $x^*$.


Далее будем предполагать, что уравнение (2)


однозначно разрешимо для любого $y \in Y$. Его решение обозначим


$x_n^*$. Возникает вопрос о нахождении среди всех однозначно разрешимых


проекционных методов тех, для которых отклонение точного решения $x^*$ от


приближенного решения $x_n^*$ было бы наименьшим на заданном множестве


решений $F \subset X$. Для характеризации этого отклонения была введена


[1]--[3] {\it оптимальная оценка погрешности класса всех


проекционных методов}


$$


V_n(F) = \inf_{X_n, Y_n, P_n} \sup_{x^*\in F} \|x^* - x_n^*\|,


$$


где \; $\inf$\; берется по всевозможным подпространствам $X_m$, $Y_m$


размерности $m \leq n$ в $X$ и $Y$ соответственно и линейным


сюръекциям $P_m$ из $Y$ на $Y_m$ таким, что уравнение (2)


однозначно разрешимо для любых $y\in Y$.





\begin{defn}


Проекционный метод (2) называется {\it


асимптотически оптимальным для множества решений $F$ среди всех


проекционных методов}, если


$$


\sup_{x^*\in F} \|x^* - x_n^*\| \sim V_n(F) \quad (n \rightarrow \infty).


$$


\end{defn}





Рассмотрим величину


$$


V(F,X_n) = \inf_{Y_n,P_n} \sup_{x^* \in F}\|x^* - x_n^*\|.


$$





\begin{defn}


Проекционный метод (2) назовем {\it асимптотически


оптимальным для множества $F$}, если


$$


\sup_{x^*\in F} \|x^* - x_n^*\| \sim V(F,X_n)\quad


(n \rightarrow \infty).


$$


\end{defn}





Говорят (напр., [5], с.\,191), что последовательность


подпространств $\{ X_n\} _{n=1}^{\infty}$ $(X_n\subset X)$


{\it предельно плотна} в $X$, если $E_n(x)\to 0$ $(n\to \infty)$


для любого $x\in X$.





Напомним (напр., [4], с.\,17), что {\it $n$-м


проекционным поперечником} центрально-сим\-мет\-рич\-ного (ц.\,с.) множества


$F\subset X$ называется величина


$$


\pi_n(F) = \inf_{X_m,S_m} \sup_{x \in F} \|x - S_mx\|,


$$


где\; $\inf$\; берется по всем $m$-мерным подпространствам $X_m$ в $X$ таким,


что $m\leq n$ и всем проекторам $S_m$ в $X$ таким, что $\im S_m = X_m$.





\begin{thm}Для любого центрально-симметричного


множества $F \subset X$ решений уравнения $(1)$ имеет место


равенство


$$


V_n(F) = \pi_n(F).


$$


\end{thm}





Для любого множества $F \subset X$ положим


$$


E(F,X_n)=\sup_{x \in F}E_n(x),


$$


где $E_n(x)$ --- наилучшее приближение элемента $x$ подпространством


$X_n$.





\begin{thm}


Пусть $X$ --- гильбертово пространство,


$K$ --- плотно заданный в $X$ инъективный оператор, $(\Ker P_n)^\bot


\subset D(K^*)$, $\dim K^*(\Ker P_n)^\bot=n$,


$B_n^1$ --- единичный шар в $K^*((\Ker P_n)^\bot)$.


Тогда, если $E(B_n^1,X_n)<1$, то уравнение $(2)$ однозначно разрешимо


для любого $y\in \im K$ и


$$


\|x^* -x^*_n\| \leq \frac{E_n(x^*)}{\sqrt{1-E^2(B_n^1,X_n)}}.


$$


В частности, если $E(B_n^1,X_n) \rightarrow 0$ $(n\rightarrow \infty)$,


то для любого ц.\,с. множества решений $F\subset X$


проекционный метод $(2)$ будет асимптотически оптимальным.


Если, кроме того, $X_n$ является экстремальным подпространством для


$\pi_n(F)$ $(n\geq n_0)$, то проекционный метод $(2)$ будет


асимптотически оптимальным среди всех проекционных методов решения


уравнения $(1)$.


\end{thm}





Пусть $H$ --- гильбертово пространство. Рассмотрим уравнение


\begin{equation}


Kx=: (A+B)x=y\quad (x,y \in H),


\end{equation}


где $A$ и $B$ --- плотно заданные в $H$ линейные операторы.





\begin{thm} Пусть существует непрерывный обратный к $A$


оператор $A^{-1}$, $B$ --- компактный оператор, последовательность


подпространств $\{ X_n\} _{n=1}^{\infty}$ предельно плотна


в $H$, $(\Ker P_n)^\bot \subset D(A^*)$ $(n \geq n_0)$ и


$E(A^*(V_n^1),X_n) \rightarrow 0$ $(n \rightarrow \infty )$, где


$V_n^1$ --- единичный шар в $(\Ker P_n)^\bot$.


Если 


однородное уравнение, соответствующее уравнению $(3)$, имеет только


нулевое решение, то проекционный метод $(2)$ решения уравнения


$(3)$ будет асимптотически оптимальным для любого ц.\,с. множества


решений $F$. Если, кроме того, $X_n$ является экстремальным


подпространством для проекционного поперечника $\pi _n(F)$ $(n\geq n_0)$,


то проекционный метод $(2)$ будет асимптотически оптимальным среди


всех проекционных методов решения уравнения $(3)$.


\end{thm}





{\bf Пример 1} представляет интегральное уравнение


\begin{equation}


\alpha(s)x(s)+ \int_{a}^{b} h(s,t)x(t)dt=y(s),


\end{equation}


где решение $x$ ищется в $L_2[a,b]$, $y \in L_2[a,b]$, $h \in


L_2[a,b]^2$, $D(\alpha)$ плотно в $L_2[a,b]$ и $1/\alpha\in L_\infty [a,b]$.





Пусть последовательность подпространств $\{ X_n\} _{n=1}^{\infty}$


предельно плотна в $L_2[a,b]$, $Y_n=X_n/a$, $P_n$ --- ортопроектор на


$Y_n$.





Предположим, что однородное уравнение, соответствующее уравнению


(4), имеет только нулевое решение.


Тогда проекционный метод (2) решения уравнения (4) является


асимптотически оптимальным для любого ц.\,с. множества решений $F$.





Если, кроме того, $X_n$ является экстремальным для проекционного


поперечника $\pi _n(F)$ $(n\geq n_0)$, то проекционный метод (2)


будет асимптотически оптимальным среди всех проекционных методов


решения уравнения (4).





\begin{thm}


Пусть существует $A^{-1}$ --- компактный оператор,


$B$ --- непрерывный оператор, последовательность подпространств


$\{ X_n\} _{n=1}^{\infty}$ предельно плотна в $H$, $(\Ker P_n)^\bot


\subset D(A^*)$ $(n \geq n_0)$ и $E(A^*(V_n^1),X_n) \rightarrow 0$


$(n \rightarrow \infty )$, где $V_n^1$ --- единичный шар в


$(\Ker P_n)^\bot$.


Если однородное уравнение,


соответствующее уравнению $(3)$, имеет только нулевое решение, то


проекционный метод $(2)$ решения уравнения $(3)$ будет


асимптотически оптимальным для любого ц.\,с. множества решений $F$. Если,


кроме того, $X_n$ является экстремальным для проекционного


поперечника $\pi _n(F)$ $(n\geq n_0)$, то проекционный метод $(2)$


будет асимптотически оптимальным среди всех проекционных методов


решения уравнения $(3)$.


\end{thm}





{\bf Пример 2} --- сингулярное интегральное уравнение


$I$-рода


\begin{equation}


\frac{1}{\pi} \int_{-1}^1 \frac{\sqrt{ 1- \tau ^2}}{\tau -t}x(\tau)d


\tau+H(x, \tau)=y(t),


\end{equation}


где $x$ ищется в $L_{2,p}[-1,1]$, $p(t)= \sqrt{1-t^2}$, $y(t)


\in L_{2,p}[-1,1]$, $H:L_{2,p} \rightarrow L_{2,p}$ ---


непрерывный оператор такой, что однородное уравнение, соответствующее


уравнению (5), имеет только нулевое решение. Пусть $X_n=


\text{Lin}\, \{1, t , \dotsc , t^{n-1} \}$,


$P_n$ --- ортопроектор на $X_n$. Тогда проекционный метод (2) решения


уравнения (5) является асимптотически оптимальным для любого ц.\,с.


множества решений $F$.


\medskip





{\bf Пример 3} составляет краевую задачу


\begin{gather*}


\frac{d^2 x}{dt^2}+p(t)x(t)=y(t),\\


x(0)=x(1)=0,


\end{gather*}


где $p(t) \in L_{\infty }[0,1]$. Пусть $X_n=\text{Lin}\,\{ \sin k\pi


t,\ k=\overline{1,n} \}$, $P_n$ --- ортопроектор на $X_n$.


Тогда проекционный метод (2) решения рассматриваемой задачи является


асимптотически оптимальным для любого ц.\,с. множества $F$.





\begin{thebibliography}{9}





\bibitem{1}


Бахвалов Н.С. {\em Численные методы. Анализ, алгебра,


обыкновенные дифференциальные уравнения}. -- М.: Наука, 1973. -- 631 c.


\bibitem{2}


Габдулхаев Б.Г. {\em Оптимальные аппроксимации решений


линейных задач}. -- Казань: Изд-во Казанск. ун-та, 1980. -- 232 с.


\bibitem{3}


Габдулхаев Б.Г. {\em Оптимальные аппроксимации решений линейных задач 


и прямые методы решения сингулярных интегральных уравнений}:


Докт. дис. в форме науч. докл. -- Киев, 1985. -- 48~c.


\bibitem{4}


Корнейчук  Н.П.


{\em Точные константы в теории приближения}. --


М.: Наука, 1987. -- 424~c.


\bibitem{5}


Красносельский М.А., Вайникко Г.М.,


Забрейко П.П., Рутицкий Я.Б., Стеценко В.Я.


{\em Приближенное решение операторных уравнений}. -- М.:


Наука, 1969. -- 455 с.





\end{thebibliography}





\vskip 0.5cm





\post{\it Казанский государственный университет}{\it Поступила}


\post{}{22.12.1998}





\label{end}


\end{document}


